\documentclass[../DoAn.tex]{subfiles}
\begin{document}

\begin{center}
    \Large{\textbf{ABSTRACT}}\\
\end{center}
\vspace{1cm}
Last-mile logistics is the most expensive stage of the supply chain, and combining trucks with drones is regarded as a promising way to shorten service time. This thesis studies a pickup scheduling and routing problem for a cooperative transport system with multiple trucks, multiple drones and time windows. Each vehicle may perform several trips, picked-up goods may not stay on board longer than a given limit, some customers can only be served by trucks, and the objective is to minimize the completion time of the last vehicle (makespan). Most existing studies consider a single truck with one or a few drones and rarely combine multiple trips with time windows, while exact methods can only solve small instances because the problem is NP-hard. The thesis therefore adopts a metaheuristic approach based on Tabu Search, which has proven effective on many vehicle routing variants. The proposed method has two phases. The first phase builds an initial solution with a greedy insertion heuristic that processes customers in earliest-due-date order. The second phase improves the solution with an Adaptive Tabu Search that uses six neighborhood operators, including an operator that relocates a whole trip, a penalized objective whose weights are adjusted dynamically by Strategic Oscillation, and a Ruin \& Recreate diversification mechanism. The thesis also implements four neighborhood generation strategies and analyzes their computational complexity. The main contributions are the formal model of the problem with a proof of NP-hardness, the Adaptive Tabu Search algorithm together with an implementation that includes an independent solution validator, and an experimental evaluation on 80 instances with 6 to 100 customers. The algorithm always finds a feasible solution, matches the reference solution on 35 of the 36 instances with at most 12 customers, and stays within about 1.9\% of the reference on average over the whole data set. A comparison over five independent runs shows that choosing operators sequentially gives slightly better solutions than choosing them at random, while the random strategy is about 9\% faster.

\end{document}
